% Generated draft: requires mathematical and source review.
\documentclass[11pt]{article}
\usepackage[margin=25mm]{geometry}
\usepackage{fontspec}
\setmainfont{Latin Modern Roman}
\usepackage{amsmath,amssymb,mathtools,unicode-math}
\setmathfont{Latin Modern Math}
\usepackage{enumitem,needspace,xurl,hyperref}
\setlength{\parindent}{0pt}
\setlength{\parskip}{5pt}
\setlength{\emergencystretch}{3em}
\newcommand{\N}{\mathbb{N}}
\newcommand{\Z}{\mathbb{Z}}
\newcommand{\Q}{\mathbb{Q}}
\newcommand{\R}{\mathbb{R}}
\newcommand{\C}{\mathbb{C}}
\newcommand{\F}{\mathbb{F}}
\newcommand{\A}{\mathbb{A}}
\newcommand{\PP}{\mathbb{P}}
\newcommand{\E}{\mathbb{E}}
\newcommand{\eps}{\varepsilon}
\newcommand{\dd}{\,\mathrm{d}}
\begin{document}
\section*{UnsolvedMath ID: 30003060}
\textbf{Source problem code:} catalogue entry
\subsection{Definitions and mathematical statement}
The supplied question leaves the ground field and the definition of higher Koszul homology implicit. The following gives a conventional field-based formulation; these conventions require confirmation against the cited sources.

Let \(k\) be a field, let \(V\) be a \(k\)-vector space, and let \(R\subseteq V\otimes_k V\). The associated connected graded quadratic algebra is
\[
A=T_k(V)/(R),\qquad A_0=k,\qquad \deg(V)=1.
\]
Its augmentation gives \(k\) the structure of a graded \(A\)-module. Say that \(A\) is Koszul if
\[
\operatorname{Tor}^{A}_p(k,k)_q=0\qquad\text{for all }p,q\geq0\text{ with }p\ne q,
\]
where \(p\) is homological degree and \(q\) is internal degree.

Put \(W_0=k\), \(W_1=V\), and, for \(p\geq2\), put
\[
W_p=\bigcap_{i=0}^{p-2}V^{\otimes i}\otimes R\otimes V^{\otimes(p-2-i)}\subseteq V^{\otimes p},
\]
with tensor products over \(k\). The Koszul chain complex with coefficients in \(A\) has terms \(C_p^K(A)=A\otimes_k W_p\) and differential given, by linear restriction from tensor expressions, by
\[
b_K(a\otimes x_1\otimes\cdots\otimes x_p)
=ax_1\otimes x_2\otimes\cdots\otimes x_p
+(-1)^p x_pa\otimes x_1\otimes\cdots\otimes x_{p-1}
\qquad(p\geq1).
\]
Write \(HK_p(A)=H_p(C_\bullet^K(A),b_K)\). In the usual higher Koszul construction, the Koszul cohomology class represented by the inclusion \(e_A:V\longrightarrow A\), \(v\longmapsto v\), acts by cap product to give a square-zero differential
\[
\partial_e:HK_p(A)\longrightarrow HK_{p-1}(A).
\]
Define
\[
HK_p^{\mathrm{hi}}(A)=
\frac{\ker(\partial_e:HK_p(A)\to HK_{p-1}(A))}
{\operatorname{im}(\partial_e:HK_{p+1}(A)\to HK_p(A))},
\qquad HK_{-1}(A)=0.
\]
Here the coefficients are \(A\), and positive degrees mean positive homological degrees, with vanishing in every internal degree.

Under these conventions, the proposed equivalence is
\[
A\text{ is Koszul}
\quad\Longleftrightarrow\quad
HK_p^{\mathrm{hi}}(A)=0\text{ for every integer }p>0.
\]
The intended scope over ground fields, any characteristic restrictions, and the coefficient and grading conventions are not specified in the supplied statement. Thus the universally quantified version over all fields and all quadratic algebras is a candidate interpretation requiring source confirmation. The supplied dated assessment describes the forward implication as proved and the converse as conjectural; that assessment is not independently verified here.
\subsection{Short English statement}
The question asks whether a quadratic algebra is Koszul exactly when its higher Koszul homology vanishes in every positive homological degree. The displayed formulation uses conventional field-based definitions and coefficients in the algebra itself; the source record leaves these conventions implicit. Its supplied dated assessment labels the problem open and identifies the converse implication as conjectural.
\subsection{Sources}
\begin{itemize}
\item Local registry: \texttt{data/unsolvedmath/problems.json}, supplied record with permanent UnsolvedMath ID \(30003060\), source problem code \texttt{catalogue entry}, title ``Koszulness from Vanishing Higher Koszul Homology.'' Collection: Oberwolfach Reports Open Problems. The record gives proposed year 2016 and names no individual proposer.
\item Supplied report citation: \emph{Hochschild Cohomology in Algebra, Geometry, and Topology} (2016), Oberwolfach Reports, pp.\ 449--506, DOI 10.4171/owr/2016/10. \url{https://doi.org/10.4171/owr/2016/10}.
\item Supplied literature reference: ``Koszul calculus.'' \url{https://arxiv.org/abs/1512.00183}.
\item Supplied literature reference: ``higher Koszul calculus.'' \url{https://www.sciencedirect.com/science/article/pii/S0021869318306239}.
\item Literature assessment dated exactly \texttt{2026-08-21}: supplied dated curation, not independently verified. It labels the problem open, reports that the Koszul-calculus paper proves the vanishing implication for Koszul quadratic algebras and records the converse as a conjecture, and says that examples support the conjecture but no general proof was located. The background explicitly requests independent verification of time-sensitive status and cited claims.
\end{itemize}
\end{document}
